% Generated draft: requires mathematical and source review.
\documentclass[11pt]{article}
\usepackage[margin=25mm]{geometry}
\usepackage{fontspec}
\setmainfont{Latin Modern Roman}
\usepackage{amsmath,amssymb,mathtools,unicode-math}
\setmathfont{Latin Modern Math}
\usepackage{enumitem,needspace,xurl,hyperref}
\setlength{\parindent}{0pt}
\setlength{\parskip}{5pt}
\setlength{\emergencystretch}{3em}
\newcommand{\N}{\mathbb{N}}
\newcommand{\Z}{\mathbb{Z}}
\newcommand{\Q}{\mathbb{Q}}
\newcommand{\R}{\mathbb{R}}
\newcommand{\C}{\mathbb{C}}
\newcommand{\F}{\mathbb{F}}
\newcommand{\A}{\mathbb{A}}
\newcommand{\PP}{\mathbb{P}}
\newcommand{\E}{\mathbb{E}}
\newcommand{\eps}{\varepsilon}
\newcommand{\dd}{\,\mathrm{d}}
\begin{document}
\section*{UnsolvedMath ID: 30002344}
\textbf{Source problem code:} catalogue entry
\subsection{Definitions and mathematical statement}
A complex threefold has complex dimension \(3\). A normal complex space is \(\Q\)-factorial if every Weil divisor has a positive integral multiple that is Cartier. Terminal singularities mean that the canonical divisor \(K_X\) is \(\Q\)-Cartier and every exceptional prime divisor over \(X\) has strictly positive discrepancy: on a resolution \(f\colon Y\to X\), the coefficients of exceptional divisors in
\[
K_Y-f^*K_X=\sum_E a(E,X)E
\]
are strictly positive.

Interpret a subvariety here as a closed complex analytic subvariety. The intended simplicity hypothesis is that there exist proper closed analytic subsets \(Z_j\subsetneq X\), indexed by \(j\in\N\), such that for every
\[
x\in X\setminus\bigcup_{j\in\N}Z_j
\]
there is no irreducible closed analytic subvariety \(V\subsetneq X\) with \(x\in V\) and \(1\leq\dim_{\C}V\leq2\). This specifies the convention for a very general point. The word \emph{proper} is necessary: the supplied wording says only ``positive-dimensional subvariety,'' which would otherwise include \(X\) itself. This intended interpretation requires confirmation against the original source.

\textbf{Question.} Let \(X\) be either a smooth compact K\"ahler threefold or a normal \(\Q\)-factorial compact K\"ahler threefold with at most terminal singularities. Assume the simplicity hypothesis just stated. Must there exist a complex torus
\[
T=\C^3/\Lambda,
\]
where \(\Lambda\subset\C^3\) is a discrete cocompact lattice of rank \(6\) over \(\Z\), and a finite group \(G\) acting holomorphically on \(T\), such that
\[
X\dashrightarrow T/G
\]
is bimeromorphic? Here \(T/G\) is the complex analytic quotient, and bimeromorphic means that the spaces have isomorphic dense open subsets whose complements are proper closed analytic subsets. The action is not assumed to be free or to preserve the origin of \(T\).

\textbf{Supplied status qualification.} The registry labels this problem partially solved. Its dated assessment reports a proved restricted terminal-threefold statement without divisors, including a finite cover by a torus, while reporting an acknowledged proof gap for the broader assertion about simple compact K\"ahler threefolds. The precise hypotheses and meaning of the restricted statement and its finite cover are not supplied, so that statement is not formulated here as a theorem. These are supplied literature assessments, not independently verified conclusions.
\subsection{Short English statement}
The question asks whether every simple compact K\"ahler threefold in either of the stated smooth or terminal \(\Q\)-factorial classes is bimeromorphic to a finite holomorphic quotient of a three-dimensional complex torus. Simplicity is interpreted as excluding positive-dimensional proper analytic subvarieties through a very general point. The supplied partial-solution assessment and its reported proof gap remain qualifications requiring review.
\subsection{Sources}
\begin{itemize}
\item Supplied local UnsolvedMath registry record, permanent ID \(30002344\), title \emph{Torus Quotients among Simple Compact K\"ahler Threefolds}, source problem code \texttt{catalogue entry}, in the Oberwolfach Reports Open Problems collection. The record gives proposed year \(2013\) and names no individual proposer. Its \texttt{corrected\_verified} label and accompanying manual-review description are supplied registry metadata, not independent source verification.
\item Supplied report citation: \emph{Complex Algebraic Geometry (2014)}, Oberwolfach Reports, DOI \texttt{10.4171/owr/2013/27}, pp.\ 1563--1627. \url{https://doi.org/10.4171/owr/2013/27}. The background attributes the problem to workshop open-problem discussions documented in Oberwolfach Reports. The supplied literature list also identifies this reference as \emph{Oberwolfach Report 2013/27}.
\item Campana--Horing--Peternell, \emph{Abundance for Kahler threefolds}, as attributed in the supplied literature list. \url{https://math.univ-cotedazur.fr/u/hoering/articles/a26-abundance.pdf}.
\item \emph{Erratum and addendum on the abundance proof}, as identified in the supplied literature list. \url{https://math.univ-cotedazur.fr/~hoering/articles/erratum-addendum-abundance.pdf}.
\item Literature assessment date: \texttt{2026-08-21}. This exact date records supplied dated curation, not an independently verified literature search. The supplied assessment reports a proved restricted terminal case without divisors, including a finite torus cover, and an acknowledged proof gap for the broader simple compact K\"ahler threefold assertion. The record expressly requests independent verification of time-sensitive status and cited claims.
\end{itemize}
\end{document}
